\section{The third level in two dimensions by separate registers}
\label{sec:registers}
Spatial coordinates need not be reserved for different logical tasks.
The construction here places the frame, ruler, computation, masks and
broadcast on separate finite registers at the same sites of
$(\Z/L\Z)^2$. Throughout this section $L=n!$, volume precedes coupling,
and the gap is the next distinct energy above the entire ground space.
The finite alphabet depends on the prescribed program and a fixed integer
$\beta$, never on $L$ or $g$.

\subsection{A shared geometry with ordinary nearest-neighbour rules}
Use the frame of Section~\ref{sec:bulk-three}, now in one $xy$ plane.
Its minimum is four; every other configuration costs at least five.
This is precisely the periodic-to-open frame of
\cite[Section~6.3, equation~(133), Lemma~55]{spectralFull}.
Its complement is an open $(L-1)$-square. The four adjacent-interior
flags are $X_0,X_1,Y_0,Y_1$. Put the unchanged five-symbol ruler of
Section~\ref{sec:ruler} on an independent register, with its horizontal
rules along $x$ and vertical rules along $y$. Write
\[
 f_g(p)=-g/p+r/p^2,\quad r=1/16,\quad
 f_* =\min(0,\min_{q\geq1}f_g(q)),\quad
 \mathcal P(g)=\{p\geq2:f_g(p)=f_*\}.
\]
For $0<g<11/480$, the minimizing pitches are at least six.
There is one minimizing pitch except at
$g_p=r(2p+1)/(p(p+1))$, where $p,p+1$ tie, and
$p=2r/g+O(1)$. The planar estimate
\eqref{eq:bulk-defect-bound} applies without summing over a third axis:
\begin{equation}\label{eq:register-defect}
 E_{\rm wall}-L^2f_*\geq
 \sum_s p_s(f_g(p_s)-f_*)+\tfrac58 V_{\rm err}.
\end{equation}
The nonnegative contribution of $W$-free columns is retained in the
underlying identity \eqref{eq:bulk-defect-identity}.
In particular a defective ruler cannot gain an extensive negative
energy by evading clean-stripe costs.

A bit $v$, constant on $y$ bonds and equal to $[\text{ruler}=V]$
off $W$, aligns the divider columns in every stripe. Require $v=1$
on $W$ sites in the frame column. Thus a zero-penalty striped ruler
has a common pitch, and the frame column is one divider column.
Every uniform pitch dividing $L$ has such a filling. At a tie,
mixed pitches fail alignment. A phase bit $I$ has unit equality
penalties on both axes. The inert phase $I=0$ restricts the ruler
to $A,D$ and fixes the auxiliary classical fields to dummy values;
it does not remove any bare ruler energy. The active phase $I=1$
enables the following rules. Frame and ruler fields remain independent.

Let $F_y,T_y$ mark respectively the first and last non-$W$ rows in
a stripe. Define the cell-first flag by the actual bond after a $V$
column. All these flags have two-site defining truth tables.
Use two independent copies of the thickness counter
\eqref{eq:bulk-thickness}. In its own horizontal interval each has
Boolean fields $b,c,f$, initialized by $b=f=0$ at $F_y$ and $c_0=1$
at the interval's first column, with internal recurrences
\[
 c_{j+1}=b_jc_j,\qquad b'_j=b_j\mathbin{\mathrm{XOR}}c_j,
 \qquad f'_j=f_j\mathbin{\mathrm{OR}}b_j.
\]
The local copy restarts after every $V$ and is dummy on $V,W$.
The global copy restarts only at $X_0$, ignores interior dividers,
and is dummy on $W$ and the frame column. Increment is checked only
between non-$W$ rows, and carry only inside its specified interval.
The frame row does not interrupt either counter.
At stripe time $t$, $b_j$ is the $j$th binary bit of $t$ and
$b_j\lor f_j=[2^j\leq t]$. At $T_y$ this selects
\[
 \ell(p)=\bitlength(p-2).
\]
Copy the top indicators into $P_{\rm cell},P_{\rm global}$ on
\emph{every} $y$ bond, including $W$ and the frame row. Set the local
mask to zero on $V$ and the global mask to zero on the frame column.
The masks mark, respectively, the first $\ell$ sites after each
divider and the first $\ell$ sites of the whole framed tape.
Common pitch makes all stripe-top demands agree. Since
$\ell(p)<p-1$ for $p\geq6$, every reader has an adjacent nonreader.

Fix $A\leq_T0'$ containing $0,1,2$ and a total computable
approximation $a(j,s)\to A(j)$. Run one initialized functional
tableau, as in Section~\ref{sec:classical}, throughout the whole frame
interior, using $x$ for tape and $y$ for time. Its work and persistent
output tracks are separate; both edge-arrival fields are blank,
and inactive fields are dummy. The machine fairly finishes stages
$a(0,s),\ldots,a(s,s)$. Its output $o_t(j)$ settles after a finite
time $t_j$. The ruler's $W$ rows and $V$ columns do not cut this
tableau. The head moves at most one site per step, so the included
times $0,\ldots,L-2$ fit in width $L-1$.

Broadcast bits $B,C,D$ live on every site, including all frame and
$W$ sites. Their ordinary bond and onsite predicates are
\begin{align*}
 B(x,y+1)&=B(x,y),&D(x,y)&=C(x-1,y),\\
 C(x,y+1)&=D(x,y),&C(x,y)&=B(x,y)\quad\text{on }W.
\end{align*}
Starting at a wall row $y_0$, these imply
$B(x,y)=b(x)$ and $C(x,y_0+t)=b(x-t)$. The next wall imposes
$b(x-p)=b(x)$. Conversely every $p$-periodic word uniquely fills
$C,D$ when $p\mid L$. This is the planar version of
Lemma~\ref{lem:bulk-broadcast}; the apparent diagonal shift is two
ordinary bonds.
Impose the sole output source at $Y_1\wedge P_{\rm global}$:
\begin{equation}\label{eq:register-source}
 B=o_{L-2}(j),\qquad j<\ell(p),
\end{equation}
and impose $B=0$ wherever $P_{\rm cell}=0$.
There is no $F_y$ condition on \eqref{eq:register-source}.
The source may lie on a wall, stripe floor, stripe top or intervening
row; the registers it tests are still present there.

\begin{lemma}[Simultaneous planar completion]\label{lem:register-completion}
For each minimum frame and aligned minimizing pitch dividing $L$,
all active auxiliary penalties can vanish. Every such filling has
the same repeated reader word $[o_{L-2}(j):j<\ell(p)]$.
\end{lemma}
\begin{proof}
Fill the functional counters, copy their top masks through all rows,
and fill the initialized tableau. Repeat its final prefix followed
by zeros with period $p$, and use the preceding formulas for $C,D$.
Every constraint vanishes. Conversely, counters and masks fix the
source positions, the functional tableau fixes their values,
periodicity copies them to every cell, and the zero mask fixes all
remaining residues. The Robinson labels enter none of these tests.
Every predicate reads one site or a single bond, with finite stored
flags checked on their defining bonds. Thus this is a fixed finite
onsite/nearest-neighbour rational interaction.
\end{proof}

The source condition is essential. For example, at $L=12,p=6$ with
walls $y=0,6$, stripe floors are $1,7$. A final source row $y=11$
would disappear if it also required $F_y$. An erroneous all-zero
broadcast would then escape comparison with a nonzero final prefix.
Condition~\eqref{eq:register-source} tests that prefix at its actual
final row. Earlier machine guesses are not retained as extra readers.

\subsection{The full spectrum, exponent and critical point}
On each stripe top activate the guarded full counter of
Lemma~\ref{lem:counter} on the $\ell(p)$ local-mask sites, with
radix $2^{\beta B}$. Its first guard uses the cell-first flag; its
last guard uses the actual bond to the next nonreader. Freeze all
exterior dynamic fields at unit cost. Every controlled edge
Laplacian remains nonnegative even in inconsistent classical blocks.
Add a qubit at each site, penalized by $|1\rangle\langle1|$ when
active. On inert $y$ bonds use
$(|01\rangle-|10\rangle)(\langle01|-\langle10|)$.
This completes rational affine coefficient matrices.

\begin{lemma}[Planar full-gap identity]\label{lem:register-gap}
For each fixed $0<g<11/480$, eventually along $L=n!$,
\begin{equation}\label{eq:register-gap}
 \Delta(L,g)=\min_{p\in\mathcal P(g)}d_p,\quad
 d_p=\Delta^+K[2^{\beta A(j)}:j<\ell(p)].
\end{equation}
For $g\leq0$ the inner gap limit is zero.
\end{lemma}
\begin{proof}
The direct-sum decomposition by classical labels is
$H=\bigoplus_\sigma(E_{\rm cl}(\sigma)\ind+K_\sigma)$ with
$K_\sigma\geq0$. Completion attains the common lower bound
$E_*=4+L^2f_*$. A nonminimum frame costs at least one. A ruler
violation costs at least $5/8$ by \eqref{eq:register-defect}.
A clean wrong pitch costs at least $L\delta(g)$, with the same
positive $\delta(g)$ as in Section~\ref{sec:bulk-three}; a $W$-free
or wholly inert ruler costs $L^2|f_*|$. These exceed $5/8$ eventually.
Every wrong auxiliary predicate or nonconstant phase contributes an
additional integer penalty. Independent bare-layer lower bounds
prevent trading these costs against a lower ruler or frame energy.

In a minimizing active block there are $(L/p)^2$ identical disjoint
counters and unit-gap exterior factors. Tensor sums with nonempty
kernels have first positive gap equal to the least factor gap,
regardless of zero-mode multiplicities. For
$L-2\geq\max_{j<\ell(p)}t_j$ the reader word is settled. Put
$k_p=\beta|A\cap[0,\ell(p))|$. The full-sector bounds are
\begin{equation}\label{eq:register-bracket}
 \frac{4^{-k_p}}{64\ell(p)^2}\leq d_p
 \leq4\pi^2 4^{-k_p}\leq\frac{\pi^2}{16}<\frac58,
 \qquad d_p\leq\frac{\pi^2}{\ell(p)^2}.
\end{equation}
Here $k_p\geq3$. All other classical blocks start above this gap.
At a tie, both pitches eventually divide $n!$; alignment leaves
their two uniform blocks and their least positive gap is the minimum
in \eqref{eq:register-gap}. We use the full counter bound, including
free paths and all head sectors, without equating it to a main orbit.

For $g\leq0$ an inert all-$A$ or all-$D$ ruler with dummy auxiliaries
and a minimum frame attains energy four. Its qubit rings have a
positive single-excitation eigenvalue
$2-2\cos(2\pi/L)$, which bounds the full gap from above and tends
to zero. This suffices without an exact formula for the ring's gap.
\end{proof}

Taking logarithms of \eqref{eq:register-bracket} gives
\begin{equation}\label{eq:register-density}
 \frac{\log(1/\Delta_\infty(g))}{\log(1/g)}
 =2\beta\rho_{\ell(p)}(A)
   +O\!\left(\frac{\log\ell(p)}{\ell(p)}\right).
\end{equation}
At a tie take the larger inverse logarithm; neighboring bitlengths
differ by at most one, so an existing density gives the same limit.
The explicit relativized mechanical-word construction in
Section~\ref{sec:bulk-three} supplies a $0'$-computable set of any
prescribed $0''$-computable density $\rho\in[0,1]$.
For completeness, use $0'$-computable rationals $u_s\to\rho$,
$N_0=1$, $N_{s+1}=(s+2)N_s$, and mechanical bits of slope $u_s$
on $[N_s,N_{s+1})$. Endpoint densities differ from $u_s$ by at most
$1/(s+2)+1/N_{s+1}$. Intermediate densities are, up to $1/n$,
convex combinations of the previous endpoint density and $u_s$.
Hence all densities converge. Set the first three bits to one and
apply the ordinary limit lemma to obtain the finite ordinary program.
For a finite $\alpha\geq0$ in $\mathsf D_3$, choose
$\beta\geq\max(1,\alpha/2)$ and $\rho=\alpha/(2\beta)$.
Equation~\eqref{eq:register-density} gives exponent $\alpha$.
The nested-thirds diagonal in Corollary~\ref{cor:bulk-beyond-sides}
also supplies planar exponents outside
$\mathsf S_2\cup\mathsf P_2$.

Even when $\alpha=0$, the last bound in
\eqref{eq:register-bracket} closes the right gap. The energy density
is zero on the left and $f_*=-g^2/(4r)+O(g^4)$ on the right, since
the frame density $4/L^2$ vanishes. Thus zero is a nonanalytic
critical point. The volume limit is eventually exact at fixed $g$;
no effective bound on the finite prefix settling time is assumed.
Each reader excitation can be placed independently at any of the
$(L/p)^2$ cells. The product-ground-vector argument of
\eqref{eq:bulk-local-weight} gives local spectral weight per area
at least $1/p^2$ at fixed $g$. Receding the frame seams and taking a
local subsequential limit gives the selected product representation
with gap $d_p$, by the same local spectral-measure argument in
Section~\ref{sec:bulk-three}. This does not classify all ground states
or establish local nondegeneracy. These observations prove
Theorem~\ref{thm:register-headline}.

\subsection{Limit depth, finite data and restricted dimension bounds}
\begin{proposition}[Finite-data and dependent-limit ceilings]
\label{prop:register-ceilings}
Algebraic local entries give uniformly computable next-distinct
finite-volume gaps. A fixed finite tuple of computable local entries
also gives computable gaps, nonuniformly in that tuple. An existing
one-size exponent therefore lies in $\mathsf D_2$.
With $k$ dependent existing limits of uniformly computable data,
and existence of every required intermediate limit, the resulting
real lies in $\mathsf D_{k+1}$.
\end{proposition}
\begin{proof}
In the algebraic case root isolation decides multiplicities and
selects the first two distinct roots. For a fixed computable real
tuple $c=(c_1,\ldots,c_m)$, let
$I_c=\{P\in\Q[X_1,\ldots,X_m]:P(c)=0\}$.
This ideal has a finite rational Gr\"obner basis. Hard-code it as
finite advice. It decides which polynomial expressions vanish;
for a nonzero expression approximation eventually decides its sign.
Real-closed-field root and multiplicity algorithms relative to this
ordered coefficient field compute each finite spectrum. Complex
entries are replaced by their real and imaginary parts. The advice
need not be recoverable from programs for arbitrary coefficient names.
This is a pointwise assertion, not a uniform degeneracy-decider.
Each existing limit of a computable rational array has a Cauchy name
in $0'$ by tail-Cauchy search. Relativize this step successively.
Positivity allows the logarithms used here. Effective moduli can
remove jumps; independent parallel limits do not add their depths.
\end{proof}

In particular, Proposition~\ref{prop:critical-ceiling} is sharp
already in two dimensions for the present volume-then-coupling
convention. Spatial dimension alone does not count these dependent
limits. It still constrains an explicitly narrower architecture.
Consider classical-forced designs whose nested boxes grow along
pairwise disjoint axis blocks, whose innermost computation is
deterministic on a finite transverse strip, and whose final quantum
reader has a fixed finite state space independent of that strip's
length. In this class attaining a proper $\mathsf D_{k+1}$ value
with $k$ nested box-size limits requires $d\geq k+1$.
Indeed, if the innermost block used only one axis, its finite
computable strip-state transition would eventually cycle. An existing
finite-reader limit must have a constant value on that cycle and is
computable by evaluating any cycle state; no real-equality decision
is needed. This removes one jump. To avoid that collapse the first
block needs two axes, and each other disjoint block needs one.
Growing readers, quantum strips and nondeterministic tableaux are
outside this statement. The planar construction above shares axes
and uses a coupling limit, so is outside that architecture.

\begin{proposition}[Elementary one-dimensional benchmarks]
For rational classical finite-range open chains, every existing
size-gap exponent is zero. Rational quantum nearest-neighbour open
chains realize size-gap exponents zero, one and two.
\end{proposition}
\begin{proof}
For the classical chain choose a common denominator $D$ of the local
energies. Distinct total energies differ by at least $1/D$. From a
ground product configuration, follow single-site changes to an
excited configuration. The first exit from the ground set changes
only boundedly many bounded local terms, so the gap is at most a
constant independent of length. Onsite projectors attain zero.
For spin-$1/2$ chains put $n=|1\rangle\langle1|$ and
$t=|01\rangle\langle10|+|10\rangle\langle01|$.
The families $\sum_jn_j$, $\sum_jt_{j,j+1}$, and
$2\sum_jn_j-\sum_jt_{j,j+1}$ have the stated exponents.
The Jordan--Wigner transform gives one-particle levels respectively
$1$, $2\cos(k\pi/(L+1))$, and
$2-2\cos(k\pi/(L+1))$. The full fermionic spectrum consists of
subset sums. For hopping alone fill all negative levels; each
excitation costs a sum of absolute nonzero levels. Its gap is
$2\sin(\pi/(2(L+1)))$ for even $L$ and
$2\sin(\pi/(L+1))$ for odd $L\geq3$. Zero modes enlarge the
kernel. The third family has positive levels and vacuum ground
state, with gap $2-2\cos(\pi/(L+1))$. These are full spectra,
without an imposed particle-number sector.
\end{proof}
These benchmarks do not determine the full quantum-chain attainment
class. General one-limit ceilings and this restricted strip argument
must not be interpreted as an unrestricted spatial lower bound.
